\subsection{F-SM-28: a compact unit and compact interval do not force \(D\in\mathrm{FP}_\infty\)}
For every nonzero commutative ring \(k\), there is an example satisfying
the full unit normalization and interval hypotheses of BR-13403--05,
with the following additional properties:
\begin{itemize}
\item the monoidal unit and the interval are finitely presented
projective objects;
\item the category has a countable set of finitely presented projective
generators, and all scalar modules of morphisms between these generators
are either \(0\) or \(k\);
\item the left ringoid module \(D\) representing \(\varphi\) is not
finitely generated, so in particular is not of type
\(\mathrm{FP}_\infty\).
\end{itemize}
The desired connective localization equivalence nevertheless holds
in this example.  The example therefore disproves the proposed
automatic-\(\mathrm{FP}_\infty\) inference, not the original
equivalence.

\paragraph{Complete construction and proof.}
Let \(\mathbb N\) denote the category with one morphism \(i\to j\)
when \(i\le j\).  Put
\[
 \mathcal M=\operatorname{Fun}
     \bigl(\mathbb N,\operatorname{Ch}_{\ge0}(k)\bigr),
 \qquad
 T=\operatorname{const}_{\mathbb N}(k[0]).
 \tag{SM28.1}
\]
This is a cocomplete \(k\)-linear abelian category.  Give it the
pointwise tensor product of chain complexes, with the usual signed
differential; its monoidal unit is \(T\).  Define
\[
 \kappa(U)=\operatorname{const}_{\mathbb N}(U[0]),\qquad
 \Lambda(X)=\operatorname*{colim}_{i\in\mathbb N}X_i,\qquad
 \varphi(X)=H_0\Lambda(X).
 \tag{SM28.2}
\]
The constant functor \(G(C)=\operatorname{const}_{\mathbb N}(C)\)
is right adjoint to \(\Lambda\), by the universal property of colimits.
Both are exact: exactness of \(G\) is pointwise, and filtered colimits
of modules are exact, so \(\Lambda\) is exact in every chain degree.
The latter familiar exactness can also be read elementwise: a kernel
or relation in a filtered colimit is witnessed at one sufficiently
large index; a finite collection of such witnesses has a common
upper index.  Thus kernels and cokernels of short exact sequences
commute with this colimit.

For a scalar module \(U\), a chain map \(X_i\to U[0]\) is precisely
a map \(H_0X_i\to U\).  Taking compatible families gives
\[
 \operatorname{Hom}_{\mathcal M}(X,\kappa U)
   \cong\operatorname{Hom}_k
       \left(\operatorname*{colim}_{i}H_0X_i,U\right)
   =\operatorname{Hom}_k(\varphi X,U).
 \tag{SM28.3}
\]
Hence \(\varphi\dashv\kappa\).
The scalar functor \(\kappa\) is exact, preserves all inverse limits,
and is compatible with the tensor products.  It also has the required
right adjoint: since \(0\) is initial in \(\mathbb N\), a map
\(T\to X\) is determined by an arbitrary element of \((X_0)_0\);
thus \(\operatorname{Hom}_{\mathcal M}(T,X)=(X_0)_0\), and
scalar tensor--Hom gives \(\kappa\dashv\operatorname{Hom}_{\mathcal M}(T,-)\).

For \(a\ge1\), let \(C_a=D^a(k)\) be the disk with identity
differential from degree \(a\) to degree \(a-1\), and set
\(C_0=k[0]\).  For each \(i,a\ge0\), define
\[
 (P_{i,a})_j=
 \begin{cases}0,&j<i,\\ C_a,&j\ge i,\end{cases}
 \tag{SM28.4}
\]
with identity transition maps after \(i\).  These free extensions
satisfy
\[
 \operatorname{Hom}_{\mathcal M}(P_{i,a},X)
   =\operatorname{Hom}_{\mathrm{Ch}}(C_a,X_i)
   =(X_i)_a.
 \tag{SM28.5}
\]
The formula for \(a=0\) is the sphere adjunction and the formula for
\(a\ge1\) is the left disk adjunction.  Evaluation is exact and
commutes with filtered colimits.  Therefore every \(P_{i,a}\) is
projective and finitely presented.  They generate: for every element
in every \((X_i)_a\), the corresponding map \(P_{i,a}\to X\) contributes
that element, and the sum of these maps is a degreewise epimorphism.
In particular
\[
 T=P_{0,0}
\]
is finitely presented and projective.

We have
\[
 \Lambda(P_{i,0})=k[0],\qquad
 \Lambda(P_{i,a})=D^a(k)\quad(a\ge1).
\]
Thus \(\Lambda\) has zero positive homology on projectives, first on
the generators and then on their coproducts and retracts.
The finite-diagonal proof BR-13401 gives
\[
 L\varphi(X)\simeq\Lambda(X).
 \tag{SM28.6}
\]
In particular, \(L\varphi(T)\to k[0]\) is the canonical identity
normalization.

Take
\[
 J=P_{0,1}=\operatorname{const}_{\mathbb N}(D^1(k)),
 \qquad
 \alpha:T\longrightarrow J
\]
to be the identity inclusion in degree zero.  The interval \(J\)
is also finitely presented and projective.
For any \(X\), at every diagram index the tensor complex has
\[
 ((J\otimes X)_i)_n=(X_i)_n\oplus(X_i)_{n-1},\qquad
 d(x,y)=(d_Xx+y,-d_Xy),
 \tag{SM28.7}
\]
where the second summand is zero for \(n=0\).
The map \(\alpha_X\) is \(x\mapsto(x,0)\), hence monic.
The maps \(h_n(x,y)=(0,x)\) satisfy \(dh+hd=1\).
They commute with the diagram structure maps, so give a natural
contraction in \(\mathcal M\).  Applying the exact colimit functor
preserves this contraction.  Equation (SM28.6) therefore gives
\(L\varphi(J\otimes X)=0\) in every degree.
The scalar compatibility (SM27.1) also holds directly.
This verifies the full interval, including its degree-zero clause.

Now take \(R\) to be the full small \(k\)-linear category of the
objects \(P_{i,a}\), with
\[
 R((i,a),(j,b))=
 \operatorname{Hom}_{\mathcal M}(P_{i,a},P_{j,b}).
\]
Restricted Yoneda identifies \(\mathcal M\) with right \(R\)-modules.
For completeness, any object has a presentation by sums of the
\(P_{i,a}\), as just proved; projectivity and finite presentation make
restricted Yoneda exact and make it preserve those sums.  Such a
presentation proves full faithfulness, and a presentation of any
right \(R\)-module by representables gives essential surjectivity
by taking the corresponding cokernel in \(\mathcal M\).

By (SM28.5), a scalar morphism module is zero when \(i<j\), and
when \(i\ge j\) it is
\(\operatorname{Hom}_{\mathrm{Ch}}(C_a,C_b)=(C_b)_a\).
For \(b=0\) this is \(k\) only when \(a=0\).
For \(b\ge1\) it is \(k\) precisely when \(a=b\) or \(a=b-1\).
Thus each scalar morphism module is \(0\) or \(k\).
All hypotheses required for the exact-adjunction construction
BR-13405 are satisfied.

The left \(R\)-module \(D\) representing \(\varphi\) is
\[
 D(i,a)=\varphi(P_{i,a})=
 \begin{cases}k,&a=0,\\0,&a\ge1.\end{cases}
 \tag{SM28.8}
\]
Furthermore
\[
 R((i,0),(j,0))=
 \begin{cases}k,&j\le i,\\0,&j>i,\end{cases}
 \tag{SM28.9}
\]
and the nonzero arrows act on (SM28.8) by their scalar coefficients.
If \(D\) were finitely generated, it would be a quotient of a finite
sum of left representables.  A map from a representable based at
\((i,a)\) to \(D\) is specified by an element of \(D(i,a)\).
For \(a\ge1\) it is zero.  For the remaining finitely many indices
\((i,0)\), choose an upper bound \(N\).  At every \((j,0)\) with
\(j>N\), all those representables evaluate to zero by (SM28.9).
Their image cannot equal \(D(j,0)=k\ne0\).
This contradiction proves that \(D\) is not finitely generated.

Finally, this construction still models the desired connective
derived category.  The counit
\(\Lambda G(C)\to C\) is an identity for every complex \(C\).
The unit \(X\to G\Lambda X\) becomes an identity after \(\Lambda\),
hence belongs to \(W=\Lambda^{-1}(\mathrm{qis})\).
Also \(G\) sends quasi-isomorphisms to \(W\), since \(\Lambda G\)
is the identity.  The ordinary adjunction consequently descends
with invertible unit and counit, and gives
\[
 W^{-1}\mathcal M\simeq D_{\ge0}(k).
\]
For an adjunction formed from any representable resolution as in
BR-13405, the same conclusion, including the raw counit criterion,
also follows from F-SM-24's finitely generated scalar-morphism case.
Thus the failure of finite generation for \(D\) creates no
counterexample to the original claim in this example. \(\square\)

\paragraph{Source scope and consequence for the search.}
A fresh visual reading of manuscript footers 47 and 48 found the
conditions (a)--(d), followed by the corollary's finitely presented
projective generators and its scalar-projectivity marginal condition.
Those lines do not state finite presentation of the monoidal unit.
The example above supplies that stronger property anyway, together
with finite presentation and projectivity of the interval itself.
Therefore neither adding a compact unit nor adding this compact
interval makes F-SM-27 an automatic solution of the full general-ring
question.  This is an explicit construction, not an inference from
an unread part of the manuscript.

